\clearpage
\SourceLeaf{609}
\ArithmeticHead{580}{L}{INHALT.}
\begingroup
\newcommand{\ContentsRow}[5]{\par\noindent\hangindent=#1\hspace*{#1}#2\hspace{.4em}\mbox{\footnotesize #3}\hfill\mbox{#4\quad #5}\par}
\ContentsRow{2em}{3.\ \textsc{Gauss} an \textsc{Drobisch}, 1834 August 14}{.\,.\,.\,.\,.\,.\,.\,.\,.\,.\,.\,.\,.\,.\,.\,.\,.}{Seite}{106}
\ContentsRow{2em}{4.\ \textsc{Gauss} an \textsc{Schumacher}, 1840 Juni 20}{.\,.\,.\,.\,.\,.\,.\,.\,.\,.\,.\,.\,.\,.\,.\,.\,.}{—}{108}
\ContentsRow{2em}{5.\ \textsc{Gauss} an \textsc{Möbius}, 1849 August 13}{.\,.\,.\,.\,.\,.\,.\,.\,.\,.\,.\,.\,.\,.\,.\,.\,.\,.}{—}{109}
\ContentsRow{2em}{Bemerkungen}{.\,.\,.\,.\,.\,.\,.\,.\,.\,.\,.\,.\,.\,.\,.\,.\,.\,.\,.\,.\,.\,.\,.\,.\,.\,.\,.}{—}{109}
\par\medskip\noindent\emph{Nachlass und Briefwechsel.}\par
\noindent\hspace*{.5em}\textls[120]{Über die Kreisteilungsgleichung.}\par
\ContentsRow{1em}{I.\quad Über die Perioden von $\dfrac{p-1}{3}$ und $\dfrac{p-1}{4}$ Gliedern}{.\,.\,.\,.\,.\,.\,.\,.\,.\,.\,.\,.\,.}{—}{111}
\ContentsRow{3em}{Bemerkung}{.\,.\,.\,.\,.\,.\,.\,.\,.\,.\,.\,.\,.\,.\,.\,.\,.\,.\,.\,.\,.\,.\,.\,.\,.}{—}{113}
\ContentsRow{1em}{II.\quad Der goldene Lehrsatz}{.\,.\,.\,.\,.\,.\,.\,.\,.\,.\,.\,.\,.\,.\,.\,.\,.\,.\,.\,.\,.\,.\,.\,.}{—}{114}
\ContentsRow{3em}{Bemerkung}{.\,.\,.\,.\,.\,.\,.\,.\,.\,.\,.\,.\,.\,.\,.\,.\,.\,.\,.\,.\,.\,.\,.\,.\,.}{—}{115}
\ContentsRow{1em}{III.\quad Die Unzerlegbarkeit der Kreisteilungsgleichung}{.\,.\,.\,.\,.\,.\,.\,.\,.\,.\,.\,.\,.\,.\,.}{—}{116}
\ContentsRow{3em}{Bemerkung}{.\,.\,.\,.\,.\,.\,.\,.\,.\,.\,.\,.\,.\,.\,.\,.\,.\,.\,.\,.\,.\,.\,.\,.\,.}{—}{119}
\ContentsRow{1em}{IV.\quad Über das regelmäßige Siebenzehneck}{.\,.\,.\,.\,.\,.\,.\,.\,.\,.\,.\,.\,.\,.\,.\,.\,.\,.\,.}{—}{120}
\ContentsRow{3em}{1.\ \textsc{Pfaff} an \textsc{Gauss}, 1802 März 22}{.\,.\,.\,.\,.\,.\,.\,.\,.\,.\,.\,.\,.\,.\,.\,.}{—}{120}
\ContentsRow{3em}{2.\ \textsc{Gauss} an \textsc{Gerling}, 1819 Januar 6}{.\,.\,.\,.\,.\,.\,.\,.\,.\,.\,.\,.\,.\,.\,.\,.}{—}{121}
\par\medskip\noindent\hspace*{.5em}\textls[120]{Allgemeines zur Lehre von den Gleichungen.}\par
\ContentsRow{1em}{V.\quad Die Beziehungen zwischen den Koeffizienten der Gleichung und den Potenzsummen ihrer Wurzeln}{.\,.\,.\,.\,.\,.\,.\,.\,.\,.\,.\,.\,.\,.\,.\,.\,.\,.\,.\,.\,.\,.\,.\,.\,.}{—}{127}
\ContentsRow{3em}{Bemerkung}{.\,.\,.\,.\,.\,.\,.\,.\,.\,.\,.\,.\,.\,.\,.\,.\,.\,.\,.\,.\,.\,.\,.\,.\,.}{—}{128}
\ContentsRow{1em}{VI.\quad Algebraischer Lehrsatz}{.\,.\,.\,.\,.\,.\,.\,.\,.\,.\,.\,.\,.\,.\,.\,.\,.\,.\,.\,.\,.\,.\,.}{—}{129}
\ContentsRow{3em}{1.\ Nachlaß}{.\,.\,.\,.\,.\,.\,.\,.\,.\,.\,.\,.\,.\,.\,.\,.\,.\,.\,.\,.\,.\,.\,.\,.\,.}{—}{129}
\ContentsRow{3em}{2.\ \textsc{Gauss} an \textsc{Schumacher}, 1833 April 2}{.\,.\,.\,.\,.\,.\,.\,.\,.\,.\,.\,.\,.\,.}{—}{130}
\ContentsRow{3em}{3.\ \textsc{Gauss} an \textsc{Schumacher}, 1836 Juni 20}{.\,.\,.\,.\,.\,.\,.\,.\,.\,.\,.\,.\,.\,.}{—}{131}
\ContentsRow{3em}{4.\ \textsc{Gauss} an \textsc{Schumacher}, 1836 Juni 24}{.\,.\,.\,.\,.\,.\,.\,.\,.\,.\,.\,.\,.\,.}{—}{132}
\ContentsRow{3em}{Bemerkungen}{.\,.\,.\,.\,.\,.\,.\,.\,.\,.\,.\,.\,.\,.\,.\,.\,.\,.\,.\,.\,.\,.\,.\,.}{—}{133}
\begin{center}\rule{.13\linewidth}{.4pt}\end{center}
\begin{center}\bfseries ANALYSIS.\end{center}
\ContentsRow{0pt}{Nachträgliche Bemerkung zu der S.\ 36—64 des VIII.\ Bandes abgedruckten Abhandlung: \emph{De integratione etc.}}{.\,.\,.\,.\,.\,.\,.\,.\,.\,.\,.\,.\,.\,.\,.\,.\,.\,.\,.\,.\,.\,.\,.\,.\,.\,.\,.\,.}{—}{137}
\par\medskip\noindent\emph{Nachlass und Briefwechsel.}\par
\ContentsRow{.5em}{Exercitationes Mathematicae}{.\,.\,.\,.\,.\,.\,.\,.\,.\,.\,.\,.\,.\,.\,.\,.\,.\,.\,.\,.\,.\,.\,.\,.}{—}{138}
\ContentsRow{3em}{Bemerkungen}{.\,.\,.\,.\,.\,.\,.\,.\,.\,.\,.\,.\,.\,.\,.\,.\,.\,.\,.\,.\,.\,.\,.\,.}{—}{143}
\par\medskip\noindent\hspace*{.5em}\textls[120]{Alteste Untersuchungen über lemniskatische Funktionen.}\par
\ContentsRow{1em}{I.\quad Reihenentwicklungen und Additionstheoreme}{.\,.\,.\,.\,.\,.\,.\,.\,.\,.\,.\,.\,.\,.\,.\,.}{—}{145}
\ContentsRow{3em}{Bemerkungen}{.\,.\,.\,.\,.\,.\,.\,.\,.\,.\,.\,.\,.\,.\,.\,.\,.\,.\,.\,.\,.\,.\,.\,.\,.}{—}{149}
\ContentsRow{1em}{II.\quad Doppelte Periodizität, Produktdarstellung und Reihenentwicklungen von Zähler und Nenner}{.\,.\,.\,.\,.\,.\,.\,.\,.\,.\,.\,.\,.\,.\,.\,.\,.\,.\,.\,.\,.\,.\,.\,.\,.\,.\,.}{—}{152}
\ContentsRow{1em}{III.\quad Teilung der Lemniskate}{.\,.\,.\,.\,.\,.\,.\,.\,.\,.\,.\,.\,.\,.\,.\,.\,.\,.\,.\,.\,.\,.\,.}{—}{160}
\ContentsRow{3em}{Bemerkungen zu den Abschnitten II und III}{.\,.\,.\,.\,.\,.\,.\,.\,.\,.\,.\,.\,.}{—}{164}
\ContentsRow{1em}{IV.\quad Vermischte Formeln zur Theorie der lemniskatischen Funktionen}{.\,.\,.\,.\,.\,.\,.\,.}{—}{167}
\ContentsRow{3em}{Bemerkungen}{.\,.\,.\,.\,.\,.\,.\,.\,.\,.\,.\,.\,.\,.\,.\,.\,.\,.\,.\,.\,.\,.\,.\,.}{—}{170}
\endgroup
